%% Chapitre 06 — Vérifier le travail de l'agent
%% RÈGLE : uniquement des commandes sémantiques bookforge. Aucun \chapter,
%% \textbf, \emph, \begin{lstlisting}, \vspace... brut.

\chapitre{Vérifier le travail de l'agent}

\introChapitre{
  Accepter un diff sans le comprendre, c'est signer un chèque en blanc. Ce
  chapitre durable est le cœur de la promesse du livre : des méthodes concrètes
  pour relire, tester et challenger ce que produit l'agent, sans devoir tout
  réécrire à la main. Vérifier vite et bien est ce qui rend la délégation
  rentable.
}

\section{Pourquoi la vérification n'est pas optionnelle}

\paragraphe{Tu viens de voir que l'agent ne raisonne que sur ce qu'il a sous
les yeux. Même parfaitement briefé, avec le bon contexte et une tâche bien
bornée, il se trompe. Pas par malice : un agent produit ce qui ressemble
statistiquement à du code correct, ce qui n'est pas la même chose que du code
correct. La plupart du temps les deux coïncident. Parfois non, et c'est là que
tu dois être présent.}

\paragraphe{Le danger porte un nom : l'\terme{hallucination}. L'agent affirme
qu'une fonction existe alors qu'elle n'existe pas, suppose un comportement
d'API qu'il n'a jamais vérifié, ou déroule un raisonnement convaincant et faux.
Le piège, c'est que ça \emphase{paraît} juste. Le code compile, se lit bien,
porte les bons noms de variables. Un humain pressé valide.}

\paragraphe{La vérification n'est donc pas une corvée de fin de tâche, c'est la
moitié du travail de délégation. \important{Si vérifier te coûte plus cher que
faire toi-même, tu n'as rien gagné.} Tout ce chapitre vise un seul objectif :
rendre ta relecture rapide et fiable, pour que le temps gagné à déléguer ne
soit pas reperdu à inspecter.}

\section{Lire un diff efficacement}

\paragraphe{Le diff est ton unité de contrôle. Ne le lis pas comme un texte,
de haut en bas : lis-le par priorité. Trois questions, dans l'ordre.}

\paragraphe{\important{Un : ai-je demandé exactement ça ?} Repère les fichiers
touchés. Si tu attendais une modification dans un module et que l'agent en a
changé trois autres, arrête-toi. Un diff plus large que la tâche est le premier
signal d'alerte, avant même de lire une seule ligne de logique.}

\paragraphe{\important{Deux : qu'est-ce qui a été supprimé ?} Les lignes
ajoutées attirent l'\oe il ; ce sont les suppressions qui cassent la
production. Un agent qui « nettoie » peut retirer une garde, un cas limite, un
\code{try/catch} qui te semblait inutile mais protégeait d'un crash réel.}

\paragraphe{\important{Trois : la logique tient-elle aux frontières ?} Concentre
ta lecture sur les bornes : entrée vide, valeur nulle, collection à zéro
élément, erreur réseau. C'est là que le code plausible se fissure.}

\begin{astuce}
Quand un diff est trop gros pour être relu sereinement, c'est un défaut de
cadrage, pas de relecture. Renvoie l'agent découper son changement en
plusieurs étapes plus petites, chacune vérifiable indépendamment. Un petit
diff bien compris vaut mieux qu'un gros diff approuvé de confiance.
\end{astuce}

\section{Faire produire et exécuter des tests}

\paragraphe{La relecture détecte ce que tu sais chercher. Les tests attrapent
le reste. Le réflexe le plus rentable du livre tient en une phrase :
\important{ne laisse pas l'agent te dire que ça marche, fais-le prouver.}}

\paragraphe{Concrètement, demande le code et les tests dans le même mouvement,
puis exige l'exécution réelle, pas une promesse.}

\begin{extraitcode}{text}
Ajoute la fonction parsePrice, puis écris des tests qui couvrent :
prix entier, prix décimal, chaîne vide, devise inconnue.
Lance la suite et colle-moi la sortie complète, pas un résumé.
\end{extraitcode}

\paragraphe{La nuance compte. Un agent à qui tu demandes « est-ce que ça
passe ? » répondra volontiers « oui, tout est vert » — parfois sans avoir rien
lancé. En exigeant la \emphase{sortie brute} du lanceur de tests, tu remplaces
une affirmation par une preuve que tu peux relire.}

\begin{avertissement}
Méfie-toi des tests écrits par l'agent qui valident sa propre erreur. Si le
code et le test partent de la même hypothèse fausse, ils s'accordent et le vert
ment. Lis au moins les cas limites des tests : un test qui n'assère rien de
gênant ne prouve rien.
\end{avertissement}

\paragraphe{Quand un bug existe déjà, inverse l'ordre : fais d'abord écrire un
test qui échoue et reproduit le bug, puis demande le correctif. Le passage du
rouge au vert est une preuve bien plus solide qu'un « c'est corrigé ».}

\section{Demander à l'agent de justifier ses choix}

\paragraphe{L'agent est un excellent relecteur de lui-même, à condition de le
forcer à expliciter. Avant d'accepter un changement non trivial, demande-lui de
défendre ses décisions. La question fait souvent remonter le doute que la
relecture seule manque.}

\begin{extraitcode}{text}
Avant que je valide : explique pourquoi tu as choisi cette approche
plutôt qu'une autre, et liste ce qui pourrait casser en production.
Cite les lignes de code ou de doc sur lesquelles tu t'appuies.
\end{extraitcode}

\paragraphe{Le détail qui change tout, c'est l'exigence de \emphase{citer ses
sources}. Un agent qui doit pointer la ligne exacte d'une API ne peut plus
inventer son comportement aussi facilement. S'il répond « d'après la
documentation, cette méthode renvoie null si... » sans pouvoir montrer où, tu
tiens ton hallucination.}

\begin{astuce}
Demander « quels cas n'as-tu pas gérés ? » est souvent plus révélateur que
« est-ce correct ? ». La première question invite l'agent à chercher ses
propres trous ; la seconde l'invite à se rassurer.
\end{astuce}

\section{Repérer le code plausible mais faux}

\paragraphe{Les hallucinations les plus dangereuses ne sont pas grossières,
elles sont subtiles. Voici les formes que tu apprendras à flairer.}

\begin{liste}
  \element{\emphase{L'API inventée.} Une méthode au nom parfaitement crédible
  — \code{user.getFullName()} — qui n'existe pas dans la bibliothèque. Le
  réflexe : vérifier toute fonction externe que tu ne connais pas par c\oe ur.}
  \element{\emphase{Le paramètre supposé.} Un appel avec un argument
  plausible mais erroné (un ordre inversé, une unité en secondes au lieu de
  millisecondes). Ça compile, ça tourne, ça donne un résultat faux.}
  \element{\emphase{Le commentaire qui ment.} Le commentaire décrit ce que le
  code \emphase{devrait} faire ; le code fait autre chose. Lis le code, pas le
  commentaire.}
  \element{\emphase{La gestion d'erreur cosmétique.} Un \code{catch} qui
  avale l'exception sans rien traiter, juste pour faire propre.}
\end{liste}

\paragraphe{Le fil rouge : plus une ligne te paraît évidemment correcte sans
que tu l'aies vérifiée, plus elle mérite ton attention. La confiance est le
terrain de jeu de l'hallucination.}

\section{Boucler la correction avec un retour précis}

\paragraphe{Quand tu repères un défaut, la qualité de ton retour détermine la
qualité du correctif. Un retour vague relance l'agent à l'aveugle ; un retour
précis le remet sur les rails du premier coup. Compare.}

\begin{extraitcode}{text}
Mauvais retour :
  « Ça ne marche pas, corrige. »

Bon retour :
  « parsePrice renvoie NaN sur l'entrée "" (chaîne vide).
    Attendu : 0. Le test test_empty_string échoue.
    Regarde la branche ligne 14, le cas vide n'est pas traité. »
\end{extraitcode}

\paragraphe{Le bon retour donne trois choses : le \emphase{symptôme} observé,
le résultat \emphase{attendu}, et un \emphase{point d'ancrage} dans le code.
Avec ça, l'agent n'a plus à deviner ton intention ; il corrige le bon endroit.}

\paragraphe{Garde la boucle courte. Si après deux ou trois allers-retours
l'agent tourne en rond — il « corrige » une chose et en casse une autre —
c'est le signe que le contexte s'est embrouillé. Plutôt que d'insister, repars
d'une session propre avec un brief resserré, ou reprends la main sur la partie
qui résiste. Acharner l'agent coûte souvent plus cher que finir soi-même les
dix dernières lignes.}

\resumeChapitre{
  \element{\important{Lire le diff par priorité} : vérifie d'abord que le périmètre est respecté, ensuite ce qui a été supprimé, enfin la logique aux frontières — entrée vide, valeur nulle, erreur réseau.}
  \element{\important{Exiger la preuve} : ne laisse pas l'agent affirmer que ça marche — demande la sortie brute des tests, pas un résumé.}
  \element{\important{Faire justifier} : demande à l'agent d'expliquer ses choix et de citer ses sources ; il ne peut plus inventer une API s'il doit pointer la ligne de doc.}
  \element{\important{Code plausible mais faux} : une API inventée, un paramètre dans le mauvais ordre, un commentaire qui ment, un \code{catch} vide — plus c'est évident, plus ça mérite vérification.}
  \element{\important{Retour précis} : donne le symptôme observé, le résultat attendu et un point d'ancrage dans le code — un retour vague relance l'agent à l'aveugle.}
}

\paragraphe{Tu sais maintenant contrôler ce que l'agent \emphase{a fait}.
Reste à cadrer, en amont, ce qu'il a le \emphase{droit} de faire : quelles
commandes il peut lancer, quels fichiers il peut toucher, ce qui exige ton
approbation. C'est l'objet du chapitre suivant, où la vérification a posteriori
laisse place aux garde-fous posés à l'avance.}
